\documentclass[9pt]{extarticle}
\usepackage[T1]{fontenc}
\usepackage[utf8]{inputenc}
\usepackage[french]{babel}
\usepackage{newpxtext,newpxmath}
\usepackage[a4paper,top=0.7cm,bottom=0.65cm,left=0.85cm,right=0.85cm]{geometry}
\usepackage{microtype}
\usepackage{xcolor}
\usepackage{titlesec}
\usepackage[most]{tcolorbox}
\usepackage{array}
\usepackage{booktabs}
\usepackage{multicol}

\definecolor{navy}{HTML}{1F3A5F}
\definecolor{grey}{HTML}{54606E}
\definecolor{green}{HTML}{2C6E49}
\definecolor{red}{HTML}{A6322E}
\definecolor{boxgray}{HTML}{F3F5F8}
\definecolor{boxred}{HTML}{FCF2F1}
\definecolor{boxedge}{HTML}{D9E0EA}

\setlength{\parindent}{0pt}
\setlength{\parskip}{1.1pt}
\setlength{\columnsep}{12pt}
\linespread{0.94}
\pagestyle{empty}

\titleformat{\section}{\color{navy}\bfseries}{}{0pt}{}
\titlespacing{\section}{0pt}{4pt}{0.5pt}

\newtcolorbox{these}{colback=boxgray,colframe=boxedge,boxrule=0.5pt,arc=2pt,
  left=6pt,right=6pt,top=2pt,bottom=2pt}
\newtcolorbox{alerte}{colback=boxred,colframe=red!55,boxrule=0.7pt,arc=2pt,
  left=6pt,right=6pt,top=3pt,bottom=3pt}
\newcommand{\code}[1]{\texttt{\footnotesize #1}}
\newcommand{\mm}[1]{\textbf{\textcolor{navy}{#1}}}
\newcommand{\ok}[1]{\textcolor{green}{\textbf{#1}}}
\newcommand{\ko}[1]{\textcolor{red}{\textbf{#1}}}

\begin{document}

\begin{center}
{\LARGE\bfseries\color{navy} Scella — Le Seuil \& Le Registre}\\[2pt]
{\small\color{grey}\itshape Les deux produits (ex-C3 et ex-C1), renommés, mesurés et ramenés à leur maturité réelle}\\[2pt]
{\footnotesize\color{grey} 26 août 2026 \; · \; one-pager \; · \; chaque chiffre est re-mesuré ce jour, chaque affirmation porte son fichier et sa ligne}
\end{center}
\vspace{1pt}

\begin{these}
\mm{Le but.} Faire qu'un appel sorti vers un modèle tiers laisse derrière lui un artefact que quelqu'un d'autre peut re-vérifier — sans clé, sans réseau, sans installer Scella — et qui écrit lui-même ce qu'il ne prouve pas.

\mm{Ce que ça résout.} Une organisation qui fait passer ses décisions par un modèle tiers ne peut, aujourd'hui, rien établir après coup : le journal appartient au fournisseur, il n'est pas scellé, et il ne survit pas au rollout. L'AI Act impose pourtant au système à haut risque de journaliser ses événements (art. 12) — et les passerelles existantes, Portkey, Helicone, LiteLLM, Cloudflare, routent, mesurent et mettent en cache sans produire aucun artefact opposable. L'écart visé n'est donc pas « router moins cher » : c'est \mm{rendre le passage re-prouvable}.

\mm{Les noms.} « C1 » et « C3 » sont des rangs d'inventaire : ils disent une place, jamais une fonction. \mm{Le Seuil} (ex-C3) est le point de passage — tout appel qui sort le franchit, et ce qui le franchit est scellé. \mm{Le Registre} (ex-C1) conserve ce que le Seuil a scellé et le traduit en dossier réglementaire. \emph{Le Seuil scelle ce qui passe ; le Registre garde ce qui a été scellé} — le routage est la porte, le reçu est le fossé.
\end{these}

\begin{multicols}{2}

\section{Le Seuil — ce qui tourne, et ce qui refuse de démarrer}
Le Seuil est une passerelle qui exécute un consensus de modèles et rend, avec la réponse, un reçu re-prouvable. Le fait qui décide de sa vendabilité est vérifié : \code{crates/efe-swarm-gateway/src/lib.rs:1881}, test \code{selfhost\_e2e\_real}, « E2E self-host RÉEL (Ollama local, keyless) » — consensus LLM complet, records persistés, re-preuve par \code{verify/:id}, \ok{sans aucune clé d'API}. La réserve est dans le même fichier : le test porte \code{\#[ignore]} (« nécessite un Ollama local »). C'est une capacité prouvée, \ko{pas une démo empaquetée}.

Le pendant de cette liberté est un refus explicite. \code{main.rs:411}, \code{frontier\_boot\_guard} : dès que la frontière payante est activée, quatre invariants doivent tenir, sinon le binaire ne démarre pas. Clé HMAC dérivable (sans elle « les reçus R4 sont \textsc{forgeables} »), authentification OIDC (sinon « egress tiers \textsc{payant} sans authentification »), store durable (sinon « les reçus R4 (la preuve) sont perdus à chaque rollout »), et la \emph{capacité} de caviarder — non le caviardage lui-même, mais le fait que son absence soit choisie et jamais subie.

La posture d'observation, jalon central pour qui met une passerelle dans un chemin critique, est \ok{livrée} : \code{capture.rs:140} nomme lui-même les termes usuels — \code{fail\_open} sur le chemin de capture, sous une posture \code{observe\_only} — un échec de persistance ne fait plus échouer la requête du client. Le troisième terme, \code{shadow} (capture concurrente, hors du chemin de la réponse), est \ko{délibérément non implémenté}, et le module le dit.

\section{Le Registre — un produit sans porte d'entrée à lui}
Le Registre dérive d'un run scellé le dossier réglementaire, et le sert à une console d'audit. Sa particularité structurelle doit être dite avant tout argument : \code{crates/audit-api} n'a \ko{aucun bloc \code{[[bin]]} et aucun \code{main.rs}} — c'est une bibliothèque Rust, « HTTP-as-pure-function ». Le seul exécutable du produit est \code{audit-console}. Le Registre n'a donc pas d'ingestion propre, et il n'en aura pas : son collecteur, \emph{c'est le Seuil}. \code{reg\_pont.rs} sert \code{POST /v1/swarm/reg/:id}, qui re-dérive le dossier à l'identique depuis un run déjà servi.

Cette jonction est volontairement lâche. \code{Cargo.toml:40}, la feature \code{reg-evidence} est optionnelle, et la ligne 2 du bloc dit le reste : \code{default = []}. \mm{Le Seuil ne sert rien de payant par défaut} — ni frontière, ni store durable, ni authentification, ni caviardage, ni registre. Tout cela s'active à la compilation. C'est une garantie forte pour un adoptant prudent, et c'est aussi pourquoi rien n'a jamais tourné en configuration complète.

\section{Le tissu commun}
Les deux fermetures partagent exactement quatre crates — \code{provenance-dag}, \code{scella-authority}, \code{scella-rfc3161}, \code{forensic-export} — mesurées par \code{cargo tree -e normal} intersecté avec les 74 membres du workspace. C'est peu, et c'est le bon signe : les produits se joignent par le sceau, pas par le code applicatif.

\begin{center}\footnotesize
\begin{tabular}{@{}l r r r@{}}
\toprule
 & crates & tests & lignes \code{.rs}\\
\midrule
\mm{Le Seuil} & 14 & 1\,329 & 74\,548\\
\mm{Le Registre} & 8 & 378 & 19\,386\\
partagés & 4 & 132 & 6\,043\\
\midrule
union & \textbf{18} / 74 & 1\,575 & 87\,891\\
\bottomrule
\end{tabular}
\end{center}

\section{Le reçu — et ce qu'il ne prouve pas}
Ce qui sépare le Seuil d'une passerelle de routage, c'est l'artefact : un graphe de provenance scellé, accompagné d'une ancre \emph{publiable} — opt-in par \code{EFE\_SWARM\_ANCHOR\_LOG}, et sans elle le binaire annonce lui-même au démarrage « ancrage \textsc{inactif} » — et re-prouvable \emph{hors} de l'écosystème qui l'a émis. La propriété qui rend cela crédible est vérifiable en une ligne de manifeste : \code{verify-mini} et \code{verify-meta-mini} ont tous deux un bloc \code{[dependencies]} \ok{vide}. Le cross-check de fidélité qui les compare octet à octet au producteur est en \code{[dev-dependencies]} — il ne rentre pas dans le binaire livré.

La limite est plus rare encore que la propriété, parce qu'elle est \emph{générée dans le livrable}. \code{forensic-export/src/readme.rs:32} écrit, dans le bundle remis à l'auditeur, le paragraphe qui explique comment le défaire : la déclaration « \ko{n'établit PAS} que le record n'a pas été fabriqué a posteriori (le déclarant et le détenteur de la clé peuvent être la même partie) », et, sans preuve d'ancrage externe, la re-preuve n'atteste que l'intégrité — jamais la correction, jamais la non-fabrication.

C'est de là que vient le bon cadrage. « Attestable » est une catégorie saturée : tout le monde le dit. \mm{Réfutable} est presque vide — un système construit pour être pris en défaut, par lui-même au démarrage (le garde de la frontière refuse), par n'importe qui ensuite sans clé ni dépendance (les deux vérifieurs), et qui écrit lui-même le mode d'emploi du défi. C'est la seule promesse que le code tient aujourd'hui sans réserve.

\section{Ce qu'il ne faut pas vendre}
\begin{alerte}
\mm{L'argument de coût n'a pas de source.} Le mémo du 19 août vendait le Seuil sur le prix (« $-27\,\%$ agrégé, jusqu'à $-90\,\%$ ») ; \ko{aucune de ces mesures n'existe dans le dépôt}, et le levier qui les aurait produites a été tué depuis : MZ (23 août) mesure le Cadrage à un meilleur point de $22{,}6\,\%$ de préservation, zéro palier sur douze à $99\,\%$, et établit que \code{must\_preserve} est un paramètre \emph{inerte}. \mm{Vendre le reçu, pas le prix.}
\end{alerte}
Trois autres réserves, toutes vérifiées ce jour. \ko{La frontière n'a jamais démarré en configuration complète} : ce n'est pas un manque de code — les deux jeux complets compilent, \emph{re-vérifiés ce jour} (\code{postgres,oidc,pii,reg-evidence}, puis le même \code{+frontier-llm} ; sortie 0) — mais une mise en service jamais faite. Le compilateur y laisse la trace de ce qu'aucune CI ne regarde : \code{reg\_pont.rs:252} calcule le nombre de fichiers du dossier réglementaire et le \ko{jette}, faute d'un champ pour le porter dans \code{RegReport}. \ko{Sauvegarde : zéro occurrence} dans les trois crates de service ; un registre de conservation longue sans plan de restauration est une promesse qu'un DPO démonte en une question. \ko{Rotation de clé : une seule occurrence}, \code{audit-console/src/server.rs:39}, et elle renvoie explicitement \emph{hors} du produit.

\section{Ce qui reste — et qui n'est pas du code}
Deux semaines de travail nommé — mise en service du profil complet, puis sauvegarde et rotation. Ce n'est pas ce qui manque. L'objectif est déclaré depuis le 19 août (fonder une startup) ; ce qui manque est d'un autre ordre : \ko{aucun interlocuteur extérieur n'est nommé}, ni prospect, ni client, ni partenaire de conception. La semaine décisive du plan de douze semaines — \emph{un prospect installe-t-il le Seuil en mode observation ?} — n'est donc jamais arrivée. Le risque est la distribution, pas le produit.

\section{Deux corrections nées de la vérification}
\mm{1.} Le compte de tests du Seuil circulait à \textbf{977}. C'était le compte de \code{\#[test]} seul : en ajoutant les \textbf{348} \code{\#[tokio::test]}, le chiffre réel est \textbf{1\,329}. Invisible sur le Registre, qui n'a \emph{aucun} test asynchrone : son chiffre tombait juste par accident. \mm{2.} \code{frontier\_boot\_guard} s'annonce en « triade » et son test s'appelle \code{..\_enforces\_the\_triad} ; il applique \textbf{quatre} invariants depuis M4 §2.2, et son propre corps écrit « le quatrième invariant ». Le nom est resté au stade d'avant.


\section{L'essayer}
\code{docker compose up} amène le Seuil, produit un run scellé et le re-prouve \emph{sans clé} : exécuté le 24 août (sortie 0, cinq étapes, troncature démasquée), \emph{non} rejoué ce jour — le démon Docker de cette machine est arrêté.

\end{multicols}
\end{document}
